\section{SA-DMoN --- Skeleton-Aware Deep Modularity Networks}
\label{sec_meth_sadmon}

The three learned grouping heads from
Section~\ref{sec_meth_learned_heads} all fail to improve on plain
$k$-means when applied to embeddings from the same GAT backbone.
This section considers a different approach.
The GAT backbone is retained -
but the grouping stage is replaced with a Deep Modularity Network
(DMoN)~\cite{tsitsulin2023dmon},
where the training objective is based on a differentiable relaxation
of Newman modularity~\cite{newman2006modularity}.

A vanilla DMoN configuration is tested first,
followed by nine skeleton-aware variants
that modify different parts of the modularity formulation.
Despite these changes, all ten configurations
plateau within a relatively narrow performance range
and remain below the $k$-means baseline.
The experiments point to a structural problem
with the modularity objective in this setting -
its null model is misaligned with a graph
whose edges are constructed from spatial proximity
between keypoints.

Together with the learned-head experiments of
Section~\ref{sec_meth_learned_heads},
these results provide further evidence
that changing the grouping head is not the main source of improvement.
Instead, the quality of the embedding produced by the network
is the binding constraint.

\subsection{Vanilla DMoN}
\label{sec_meth_dmon_pose}

Given the keypoint graph $\mathcal{G}$ from
Section~\ref{sec_meth_problem_graph} - with adjacency matrix
$A \in \{0,1\}^{N \times N}$ -
DMoN predicts a soft cluster-assignment matrix
$\mathbf{S} \in [0,1]^{N \times K}$.
Each row of $\mathbf{S}$ sums to one,
so that $S_{ik}$ represents the soft assignment
of node $i$ to cluster $k$.
The assignments are produced by applying a single linear projection
to the GAT embeddings followed by a softmax over the $K$ clusters.

The standard DMoN objective contains three
terms~\cite{tsitsulin2023dmon}:
\begin{equation}
    \mathcal{L}_{\text{DMoN}}
    = -\frac{1}{2m}\mathrm{Tr}\bigl(\mathbf{S}^\top \mathbf{B}\,\mathbf{S}\bigr)
      + \lambda_{\text{ortho}}\,
      \Bigl\|\mathbf{S}^\top \mathbf{S}/\|\mathbf{S}^\top \mathbf{S}\|_F - \mathbf{I}/\sqrt{K}\Bigr\|_F
      + \lambda_{\text{cluster}}\,\Bigl(\tfrac{\sqrt{K}}{N}\,\|\mathbf{s}_{\text{vol}}\|_2 - 1\Bigr),
    \label{eq:dmon_loss}
\end{equation}
where
\begin{equation}
    \mathbf{B}
    =
    \mathbf{A}
    -
    \frac{\mathbf{d}\mathbf{d}^{\top}}{2m}
\end{equation}
is the modularity matrix,
$\mathbf{d}$ contains the node degrees,
and $m = \tfrac{1}{2}\sum_{ij} A_{ij}$ is the total number of edges.
The first term encourages partitions with high modularity.
The orthogonality term discourages highly unbalanced assignments,
while the term weighted by $\lambda_{\text{cluster}}$
helps prevent solutions in which one or more clusters collapse
and receive no nodes.

Two further terms are added to the DMoN objective.
The first is a supervised cross-entropy term,
computed against the ground-truth person labels $\mathbf{y}^\star$
after Hungarian matching of clusters to people,
with weight $\lambda_{\text{supervised}}=1$.
This keeps the predicted clusters aligned with person identity,
rather than allowing the network to follow only the partition
preferred by the modularity objective.
The second is a type-exclusivity term with weight $\lambda_{\text{type}}=1$,
which penalises a cluster that accumulates
more than one joint of the same type.
In the implementation,
the spectral term itself also carries a weight,
with $\lambda_{\text{spectral}}=0.1$ unless a configuration below changes it.

This vanilla configuration forms the starting point
for the SA-DMoN experiments reported in Chapter~\ref{chap_4}.
Its main limitation comes from the structure of the graph itself.
Modularity rewards groups of nodes
that are more densely connected than expected under its null model,
while the keypoint graph is constructed from spatial proximity.
As a result, nearby joints are encouraged to remain together
even when they belong to different people.
Pose grouping often requires exactly the opposite behaviour.
This mismatch motivates the skeleton-aware modifications
to the null model introduced in the next subsection.

\subsection{SA-DMoN v1: a skeleton-aware null model}
\label{sec_meth_sadmon_v1}

The v1 variant replaces the standard degree-based null model
$\mathbf{N}_{\text{deg}} = \mathbf{d}\mathbf{d}^\top/(2m)$ with one
that includes both spatial proximity and the known structure of the
COCO skeleton.
The spatial component is defined using a Gaussian kernel
over the normalised keypoint positions
\begin{equation}
    P_{ij}
    =
    \exp\!\Bigl(
        -\,\frac{\|\mathbf{p}_i-\mathbf{p}_j\|_2^2}{2\sigma^2}
    \Bigr),
    \label{eq:sadmon_spatial}
\end{equation}
where $\sigma$ controls the spatial bandwidth
and is treated as a hyperparameter below.

The skeleton component is represented by a fixed $17 \times 17$
joint-type affinity matrix $T$,
constructed using breadth-first search over the COCO skeleton graph.
Joint types that are directly connected in the skeleton
are assigned an affinity of $T_{ab}=1.0$.
Pairs separated by two hops receive $0.5$,
and three-hop pairs receive $0.25$.
Joint types more than three hops apart -
or otherwise topologically disconnected -
receive an affinity of zero.

For nodes $i$ and $j$ with joint types $a$ and $b$, respectively,
the spatial and skeleton terms are combined as
\begin{equation}
    R_{ij}
    =
    T_{ab} \cdot P_{ij}.
    \label{eq:sadmon_null}
\end{equation}
The resulting matrix is rescaled so that its total mass matches the
edge mass of the original graph,
\begin{equation}
    \tilde{R}_{ij}
    =
    R_{ij}
    \left(
        \frac{2m}{\sum_{ij}R_{ij}}
    \right).
\end{equation}
The pose-specific modularity matrix is then
\begin{equation}
    \mathbf{B}_{\text{pose}}
    =
    \mathbf{A} - \tilde{\mathbf{R}},
\end{equation}
which replaces $\mathbf{B}$ in the DMoN loss of
Equation~\eqref{eq:dmon_loss}.
The remaining parts of the model are left unchanged,
including the assignment head,
the orthogonality and collapse-prevention terms,
and the supervised cross-entropy and type-exclusivity terms.

The bandwidth $\sigma$ proved to be the main sensitive parameter in
this formulation.
Several versions were therefore tested to determine
whether the behaviour of the null model was mainly caused by how this
bandwidth was selected.
Each configuration below corresponds to one of
the SA-DMoN ablations reported in Chapter~\ref{chap_4}.

\begin{compactenum}
    \item \textbf{Learnable, unconstrained.}
        The bandwidth is parameterised as $\sigma=\exp(\xi)$,
        with $\xi$ initialised to $\log 0.2$,
        and optimised with the rest of the network.
        This allows the model to choose its own spatial
        scale.

    \item \textbf{Learnable, clamped.}
        The same learnable parameterisation is used -
        but $\sigma$ is clamped to the range $[0.05,0.5]$
        during the forward pass.
        This prevents the bandwidth from diverging,
        as observed in the unconstrained configuration.

    \item \textbf{Learnable, clamped, $\lambda_{\text{spectral}}=1$.}
        This uses the clamped bandwidth while increasing the weight of
        the spectral term so that it has the same weighting as the
        supervised loss.
        The purpose is to test whether the modularity
        signal becomes useful when it has a larger influence on
        training.

    \item \textbf{Fixed at the per-graph median pairwise distance.}
        Here, $\sigma$ is recalculated for each scene as the median
        pairwise keypoint distance using \texttt{torch.no\_grad}.
        This gives a bandwidth that adapts to the spatial scale
        of the scene without allowing the optimiser to change it directly.

    \item \textbf{Spectral term removed.}
        The spectral weight is set to $\lambda_{\text{spectral}}=0$.
        The skeleton-aware null model is still calculated
        for diagnostic purposes -
        but it contributes no modularity gradient during training.
        This configuration separates the effect of the supervised loss
        from the modularity-based objective.
\end{compactenum}

These five v1 configurations form the methodological basis
of SA-DMoN Experiments~4--8 reported in Chapter~\ref{chap_4}.

\subsection{SA-DMoN v2: decoupled adjacency and entropy type loss}
\label{sec_meth_sadmon_v2}

The motivation for the v2 variants comes from a limitation observed in
the v1 formulation.
The spatial adjacency $\mathbf{A}$ and the spatial
component of the null model $\mathbf{P}$ are both strongly determined
by keypoint distance.
As a result, the two terms in
$\mathbf{B}_{\text{pose}} = \mathbf{A} - \tilde{\mathbf{R}}$ can
partially cancel around the selected bandwidth,
weakening the modularity signal available during training.

The v2 formulation tests two changes intended to address this problem.
The first separates the adjacency used by the modularity loss
from the geometric graph used for message passing.
The second replaces the threshold-based type-exclusivity loss
used in v1 with a smoother entropy-based objective.
The two changes are independent and can also be used together,
giving four configurations in total.

\textbf{Fix~1: feature-based adjacency:}
Instead of using the geometric $k$NN adjacency from
Section~\ref{sec_meth_problem_graph} inside the spectral loss,
a new adjacency is constructed from the current GAT embeddings.
This feature graph is formed using $k$-nearest neighbours
in embedding space with $k=8$.
The original geometric graph is still used by the GAT
for message passing -
only the adjacency supplied to the modularity objective is changed.
Since the embedding itself changes during training,
the feature-based graph is reconstructed on each forward pass.

This separates the two roles played by the graph.
Message passing still uses spatial proximity to define
which keypoints exchange information,
while the modularity objective operates on neighbourhoods
learned in the embedding space,
rather than another distance-based version of the same spatial graph.

\textbf{Fix~2: entropy type loss:}
The second change replaces the v1 type-exclusivity term
with a soft entropy-based objective
\begin{equation}
    \mathcal{L}_{\text{type}}^{\text{v2}}
    =
    -\sum_{i,k}
    S_{ik}\,\log p^{\text{type}}_{ik},
\end{equation}
where $p^{\text{type}}_{ik}$ is the estimated probability
that cluster $k$ contains the joint type $t_i$,
calculated from the current soft assignment statistics.
Compared with the threshold-based ReLU formulation used in v1,
this gives a smoother gradient as the assignments change.
The trade-off is that the sharper discrete hinge
of the earlier loss is no longer present.

\textbf{Configurations:}
Four v2 configurations are evaluated in Chapter~\ref{chap_4}.
The first uses feature-based adjacency with
$\lambda_{\text{spectral}}=1.0$.
The second uses the same feature adjacency
but reduces the spectral weight to
$\lambda_{\text{spectral}}=0.1$ -
matching the weighting used in v1.
The third keeps the original spatial adjacency
and changes only the supervised term to the entropy type loss.
The final configuration combines both changes -
using feature-based adjacency together with the
entropy-based type objective.

\subsection{Training schedule}
\label{sec_meth_sadmon_training}

All ten SA-DMoN variants are trained jointly with the GAT on the
synthetic training split.
These include the vanilla DMoN baseline,
the five v1 configurations described above,
including the no-spectral diagnostic,
and the four v2 configurations.
Training otherwise follows the optimiser settings given in
Section~\ref{sec_meth_baseline_training}.

Depending on the behaviour of the head loss,
each configuration is trained for either 100 or 150 epochs.
A 100-epoch schedule is used when the loss has clearly plateaued
by that point.
For configurations where the head loss is still changing at 100 epochs,
training is extended to 150 epochs
to allow additional time for convergence.
This avoids discarding a variant simply because its grouping objective
converges more slowly than the baseline contrastive model.

Despite the longer training schedules,
none of the SA-DMoN configurations are able to improve
improve on $k$-means applied to the baseline GAT embeddings.
The experimental results are reported in
Section~\ref{sec_res_sadmon}.
The broader reason for this behaviour -
particularly the mismatch between the modularity objective
and the structure of the pose-grouping graph -
is discussed in Section~\ref{sec_disc_modularity_autonomy_gap}.
